\documentclass[UTF8,11pt]{ctexart}
\usepackage[margin=25mm]{geometry}
\usepackage{amsmath,amssymb,booktabs,graphicx,hyperref}
\newcommand{\tabfile}[1]{\begin{center}\resizebox{\textwidth}{!}{\input{../shared/#1}}\end{center}}
\title{约束非线性动力学 SDE：\\受控可行性与尚无定论的开发集比较}
\author{PIRC-26 项目研究记录}
\date{2026 年 10 月 14 日——技术报告草稿，不代表同行评审接收}
\begin{document}
\maketitle
\begin{abstract}
本报告把已知教师恢复、真实开发集预测、分布目标与梯度路线可行性分开评价。两个基函数族均完成注册合成实验，但各自仅有 12/20 个教师匹配设置通过不变的联合恢复标准。七候选真实比较中，每个模型的 282 个声明单元格仅有 63 个可测；完整原始单元的资格样本仍为一。各预测时域的 6/6/5 单元结果是已暴露、以覆盖为条件的探索，不是合格的优势检验。非线性模型在覆盖子集中的平均 Energy 分数较低，但校准很差，60 秒平均终点误差更高。六个自编路径目标拟合及十二个配对梯度检查通过小规模可行性标准，但不证明优化收敛、一般速度或内存优势。因此，结果仅支持有限实现可行性，不支持全局预测优胜或潜状态方法资格。
\end{abstract}

\section{研究问题与结论边界}\label{sec:questions}
RQ1 研究固定设计下约束基残差能否恢复已知离散教师；RQ2 研究既有非线性拟合相对仿射比较模型是否改善注册开发总体预测；RQ3 研究轨迹分布目标及替代梯度路线能否在明确的小算例资源范围内正确运行。工程测试和合成可行性不能赋予真实数据排名资格。没有读取最终测试集，也不作等效、正式 NO\_GAIN、总体优势或部署就绪结论。

\section{模型与目标}\label{sec:model}
二维位置 $p$ 与速度 $v$ 的观测状态模型为
\begin{equation}\label{eq:model}
 dp_t=v_t\,dt,\qquad dv_t=[A v_t+B p_t+a+r_\theta(p_t,v_t,c_t)]\,dt+L\,dW_t.
\end{equation}
自编算例无上下文 $c$。位置方程固定，常数扩散在四维状态中秩为二。M0 取残差零；M1-S 使用三次样条，M1-R 使用高斯径向基；M2 使用两层宽度 16、带谱界和振幅界的神经残差。真实拟合及归一化器直接复用，不因评估结果重训。自编神经算例使用恒等归一化、单位谱界、单位振幅和残差比例一，不能据此声称重新选择了真实模型设置。

O1 在扣除仿射项后估计增量加速度。已知仿射项与常数可逆速度噪声下，合成单步响应为
\begin{equation}\label{eq:noise}
 Y_i=r_\theta(z_i)+L\xi_i/\sqrt{\Delta t},\quad\xi_i\sim N(0,I).
\end{equation}
满秩、零惩罚最小二乘协方差正比于 $(\Phi^\top\Phi)^{-1}\otimes LL^\top/\Delta t$，因子顺序依向量化约定而定。固定设计重复 $N$ 次的尺度为 $1/(N\Delta t)$。固定 $N$ 缩小步长并不必然增加漂移信息。这是事后解析解释，不是新检验、标准修改或失败的唯一因果解释。

O2 在固定训练起点及预测时域上平均分布评分。对 $K>1$ 个位置样本及目标 $y$，实际无偏配对估计量为
\begin{equation}\label{eq:energy}
 \widehat{\mathrm{ES}}=K^{-1}\sum_k\|X_k-y\|_2
 -[2K(K-1)]^{-1}\sum_{k\ne l}\|X_k-X_l\|_2.
\end{equation}
Energy 分数评价整个预测分布\cite{gneiting2007}，但估计式本身不证明校准或模型适用。终点距离使用样本均值位置，不选择最优路径。

\section{数据、设计与冻结控制}\label{sec:design}
真实冻结仍在 DSDE：\texttt{pirc20-r1t-nex326-midpoint-20260912-v1}，特征版本为 \texttt{pirc21-production-features-20260917-v2}。数据清单 SHA-256 为 \texttt{94efafdb09a62719ffbbebec07564eff640bf1adfc8ab6b1664173dc4f71abed}。本报告不分发该数据，不创建公开副本，不声称原始数据具有公开许可。保留全部七个浙江候选：六个可转换来源及一个被地理转换规则拒绝的来源。93 个保留片段各有三个时域，另有三个来源不可用单元格。没有替换候选、增加真实训练、搜索容量、插值真值或访问最终测试。

真实 M0/M2 使用原四训练单元拟合、因果起点、原生观测真值、128 条共同根路径、CPU float64 和一个数值线程。时域为 10/30/60 秒，原生结点采用时域/1024 名义分辨率，另以两倍细网格作配对诊断。共同路径减少配对模拟噪声，但重复路径和片段不是独立实验单位。新真实分析源码冻结于 Git \texttt{9ce90059f3611089150928d47aaa0393b699a016}；较早 M1 的源码时期在证据目录单独绑定，不重绑定为后续工程合并版本。

M1 含两种已知残差教师、五个固定噪声种子、0.02/0.1 秒步长、64/256 条路径及 65 个固定速度起点：每族 40 次拟合，其中 20 个教师匹配设置。已知仿射项与扩散冻结，仅改变一个速度特征。开放三次样条结点、七个 RBF 中心、系数及神经教师控制均见复现配方。真值来自离散 Euler 教师，不是连续 SDE 精确解。资格要求全部 20 个匹配设置同时满足函数 RMSE 不超过 0.25、系数相对 $\ell_2$ 误差不超过 0.5。嵌套路径数和共同噪声不构成独立重复。

\section{RQ1：恢复未取得资格}\label{sec:rq1}
\tabfile{m1-recovery.tex}
两族各完成 40 次拟合及 30 次注册滚动预测，注册路径全部有效；各仅通过 12/20 个匹配设置，均保留 \texttt{NOT\_PASSED}。最强子组五个种子全部通过不能覆盖全设置标准。表中最大设计条件数是诊断，不证明条件性是系数恢复失败的唯一原因。M1-S/M1-R 仅拟合壁钟合计约 12.59/9.86 秒，不含整作业准备、记账及等待，不能视为全部研究成本。关闭残差的比较是同一已知仿射回退，不是独立训练的公平真实 M0 基准。

\section{RQ2：已暴露且不完整的开发比较}\label{sec:rq2}
每模型均有 282 个声明单元格：63 个可测、216 个无因果原生覆盖、三个来源不可用，数值失败及未执行均为零。缺失不计为零误差。原完整单元资格规则仍只有 $n=1$；部分覆盖产生时域条件总体 6/6/5，不等于完整单元资格样本。先在原注册来源单元内平均片段，再对单元等权；本报告没有重新证实来源单元间独立性。
\tabfile{real-energy.tex}
$n$ 为覆盖原始单元数，$C$ 为配对可测片段数，$\Delta=\mathrm{M2}-\mathrm{M0}$。区间采用 2,000 次原单元配对百分位 bootstrap，条件为有因果覆盖且两模型均可测。结果已暴露，非盲确认区间，未校正多重比较，也不能修复选择或缺失；30 秒区间不跨零不能证明全局优胜。
\tabfile{fde-mc.tex}
非线性模型的 60 秒终点误差更差，即使该子集 Energy 分数更低。Point SE 是逐点 Energy jackknife 标准误的平均值，\emph{不是}单元聚合分数或其区间的标准误。受控诊断中的速度误差以观测后向差分为代理，不是潜在速度真值。
\tabfile{real-calibration.tex}
$c_x,c_y$ 为非起点原生点落入名义 95\% 坐标边际区间的等单元加权比例，既不是联合最高密度覆盖，也不是独立 Bernoulli 试验。两模型覆盖均很差。RMS 为配对粗细网格终点差异（米），非认证误差界。耦合、较小 Monte Carlo 标准误或路径无失败不能证明模型适用。每模型各时域请求路径实例为 6,656/4,352/5,120，失败零，但这些不是独立实验主体。

\section{RQ3：有限目标与梯度可行性}\label{sec:rq3}
\subsection{G1：直接路径目标拟合}
三个自编种子 26101411--26101413 与 0.25/0.125 秒步长组成六个单元格。每种子的 O1 准备使用 512 个单步转移、实际 40 次更新。师生共享已记录的约束架构，教师不从真实数据拟合。每次 O2 拟合用八条路径、两个训练起点、10/30 秒时域、最多十二次更新、学习率 0.005、梯度裁剪 100 和六步时域课程。另两个自编挑战起点不参与选检查点。恢复的是固定、独立 Brownian 训练监视器在全训练起点及全时域上的最佳检查点，不一定是最后一次迭代。
\tabfile{o2-feasibility.tex}
六个单元格全部通过预设小例标准：路径有效，四个中心差分与自动微分一致（扰动 $10^{-5}$，绝对容差 $10^{-5}$，相对容差 $10^{-3}$），仅拟合壁钟不超过 600 秒，进程生命周期峰值 RSS 不超过 2 GiB。全部以 \texttt{MAX\_STEPS} 结束，不代表收敛。600 秒/2 GiB 是测量资格门槛，不是新增研究终止硬上限。安全摘要保留全部 24 对自编挑战分数，既有改善也有恶化。训练网格没有逐路径耦合，因此跨网格结果不是隔离离散误差实验或收敛证明。O1/O2 目标不同，不直接排名其损失。

\subsection{G2/G3：配对前向问题}
三个种子、两个网格均使用同一固定 O1 准备检查点，不作模型更新。G2 配对直接 reversible Heun 与其可逆伴随；G3 配对直接 midpoint 与连续伴随 midpoint。路线参考随机伴随及可逆求解器工作\cite{li2020,kidger2021}，但本报告只评价实际安装实现与十二个算例。同一对内前向求解器一致，跨求解器差异不能当作仅梯度方法的效应。常数状态无关扩散使算例中的 It\^o/Stratonovich 漂移相同。

路线使用同一 halfway-tree Brownian 构造、每种子固定熵、$10^{-9}$ 容差，并实测粗增量等于细增量之和。每对有三个全新顺序 helper：差分验证、直接路线及候选路线；基准顺序平衡。全梯度比较覆盖所有非空加速度参数，排除零大小上下文张量；每对检查四个差分坐标。G2 全梯度相对 $\ell_2$ 容差为 $10^{-3}$，G3 为 0.05，另有配方中固定的坐标、前向及资源标准。
\tabfile{gradient-routes.tex}
十二对全部通过，配对前向差异均为零；最大相对梯度差约为 $2.801\times10^{-14}$（G2）、$3.323\times10^{-4}$（G3）。这些小 CPU 算例中，每个候选都比配对直接路线慢。RSS 是含导入与元数据的全新进程生命周期峰值，不是阶段隔离工作内存。不声称一般内存优势、渐近比较、训练后 O2 结果或真实数据加速。

\section{潜状态及矩阵方法处置}\label{sec:latent}
当前约定有观测状态，但未提供观测似然、噪声校准、潜初始分布、识别模型或可识别性条件。因此潜状态估计未获资格、未实测，不能推断其一般无必要。对 $G=[0;L]$，合法 Girsanov 漂移变化须落在 $\mathrm{range}(G)$ 并保留 $dp=v\,dt$；伪造位置噪声或加逆无效。一般高斯边际流可能违反固定位置漂移，不能直接替换为合法路径分布构造。这仅反驳该替换，不反驳所有潜流方法\cite{bartosh2025}。

瞬时秩二不等于有限时间转移奇异。对仿射 $F=[0,I;B,A]$ 及可逆 $L$，可控性块
\[[G,FG]=\begin{bmatrix}0&L\\L&AL\end{bmatrix}\]
秩为四，正时间协方差 Gramian 可非奇异。精确仿射联合转移仍合法，不能与奇异四维 Euler 协方差求逆混淆。当前约定下 L1/L2 潜组件及 G4 矩阵梯度路线保持 \texttt{NOT\_ACTIVATED}，无虚构 ELBO、潜匹配或 G4 基准值。G3 不需要扩散求逆，也不是未执行潜实验的替代标签。

\section{成本、失败与复现}\label{sec:repro}
\tabfile{costs.tex}
表为紧凑证据发布并结算时的累计计费槽秒快照。每个真实模型/方法臂保留 86,400 秒硬上限，准备、重试及孤儿结算均不重置。M0 承担共享准备及派生分析，含历史未知成本结算，因此累计数不能公平比较模型训练速度。论文构建及后续验证另有后续收据，不暗入历史快照。7,200 秒作业分配是续算/报告窗口，不是此前错误的 120 秒硬停。

早期内容写争用中断及保守费用保留。同政策重试只重算派生分析，不重复成功预测。O2 首次消费因 artifact purpose 被拒发生在数值预算分配之前；API purpose 修复改变消费授权，不改变模型、数值后缀、阈值或种子。冻结前代码可迭代，每次实际运行均保留源码、协议、规格及内容哈希；工程历史不是科学重复。

发行包包含审查后的聚合摘要、确定性表格、主张/来源索引、明确配方和验证测试。不分发完整来源、路径、单元/片段标识、轨迹、逐行预测、模型检查点、凭据或私有位置。目录分别绑定旧 M1 与新分析源码时期。公开 smoke 复现只重建表格及报告，\emph{不是}独立重做全部受控数值实验。完整科学复现需要受控访问原授权产物与归档算子。技术披露审查不等于人类科学验收，也不授权公开原数据/模型。工程 CI 或 PR 合并同样不能替代研究验收。

\section{结论}\label{sec:conclusion}
RQ1 未通过不变的全设置匹配恢复标准。RQ2 因完整单元总体资格不足、已暴露的部分比较以覆盖为条件、校准很差而尚无定论。RQ3 只确立六个直接目标及十二个配对梯度路线的小算例可行性，不证明收敛、普遍伴随优势或潜状态结果。可支持的是保留缺失与失败历史的有限负面/可行性报告，不是全局运动预测优胜模型。

\begin{thebibliography}{9}
\bibitem{gneiting2007} T. Gneiting and A. E. Raftery. Strictly proper scoring rules, prediction, and estimation. \emph{JASA} 102(477):359--378, 2007. \href{https://doi.org/10.1198/016214506000001437}{doi:10.1198/016214506000001437}.
\bibitem{li2020} X. Li, T.-K. L. Wong, R. T. Q. Chen, and D. Duvenaud. Scalable gradients for stochastic differential equations. \emph{AISTATS}, PMLR 108:3870--3882, 2020. \url{https://proceedings.mlr.press/v108/li20i.html}.
\bibitem{kidger2021} P. Kidger, J. Foster, X. C. Li, and T. Lyons. Efficient and accurate gradients for neural SDEs. \emph{NeurIPS} 34, 2021. \url{https://proceedings.neurips.cc/paper/2021/hash/9ba196c7a6e89eafd0954de80fc1b224-Abstract.html}.
\bibitem{bartosh2025} G. Bartosh, D. Vetrov, and C. A. Naesseth. Neural SDEs made easy: SDE matching. \emph{ICML}, PMLR 267:3054--3070, 2025. \url{https://proceedings.mlr.press/v267/bartosh25a.html}.
\end{thebibliography}
\end{document}
